\documentclass[10pt,a4paper]{amsart}
\usepackage[margin=19mm]{geometry}
\usepackage{amsmath,amssymb,mathtools}
\usepackage[scr=rsfs,cal=euler]{mathalfa}
\usepackage{xfrac}
\usepackage[unicode,hidelinks,bookmarks=false]{hyperref}
\newcommand{\ie}{i.e.\ }
\newcommand{\cf}{cf.\ }
\newcommand{\resp}{respectively\ }
\newcommand{\Subsection}{\subsection}
\providecommand{\nobreakdash}{\nobreak\hbox{-}}
\setlength{\emergencystretch}{2em}
\multlinegap0pt
\usepackage{mathtools}
\usepackage[shortlabels]{enumitem}
\usepackage{csquotes}
\usepackage{float}
\usepackage{amssymb}
\usepackage{xcolor}
\newtheorem{proposition}{Proposition}
\newtheorem{lemma}{Lemma}
\newcommand{\E}{\mathbb{E}}
\newcommand{\N}{\mathbb{N}}
\newcommand{\PP}{\mathbb{P}}
\DeclareMathOperator{\cD}{\mathcal{D}}
\DeclareMathOperator{\cE}{\mathcal{E}}
\newcommand{\eps}{\varepsilon}
\newcommand{\lhs}{{l.h.s.~}}
\newcommand{\one}{\textup{\textbf{1}}}
\newcommand{\op}{o_{\p}}
\newcommand{\p}{\mathbb{P}}
\newcommand{\pan}{\p^{\textup{an}}}
\newcommand{\quenchG}{\mathbf{P}^G}
\newcommand{\rhs}{{r.h.s.~}}
\newcommand{\s}{{s_\eps}}
\newcommand{\semi}{{Q}}
\newcommand{\tent}{t_{\textup{ent}}}
\newcommand{\trex}{\textit{Tx}}
\newcommand{\veps}{V^{\star}_\eps}
\newcommand{\whp}{{w.h.p.~}}
\title{Mixing trichotomy for random walks\\on directed stochastic block models}
\author{Alessandra Bianchi, Giacomo Passuello, Matteo Quattropani}
\date{Source: arXiv:2504.06851v2 (CC BY 4.0)}
\begin{document}
\maketitle

\noindent\emph{Source and licence.} Mixing trichotomy for random walks\\on directed stochastic block models. Version arXiv:2504.06851v2 (CC BY 4.0), distributed by arXiv under the Creative Commons Attribution 4.0 International licence, http://creativecommons.org/licenses/by/4.0/, as recorded in the arXiv licence metadata for this exact version and shown on the abstract page https://arxiv.org/abs/2504.06851v2. Full licence notice: \texttt{licenses/arxiv-2504.06851-CC-BY-4.0.txt}.\par\smallskip

\noindent\textit{Editorial scope.} Counted unit: the article's proposition prop:quenched-law (source0.tex lines 886--894). The article's complete original proof of it is delivered with it (source0.tex lines 895--983, one \texttt{proof} environment of 89 lines). No mathematics is written, repaired or re-derived: the statement and the proof are the article's own lines, in the article's own notation, and no retained line is substituted or re-flowed.  Retained uncounted as the source-local context the proof needs: lines 717--722; lines 750--762; lines 798--808; lines 804--806.  Nothing was omitted from any retained span, so the package carries no omission record.  No retained line contains a citation, so the wrapper carries no bibliography and no bibliography entry.  No retained line contains a figure environment, an image-inclusion command or drawing code, so no image is delivered and the package declares no asset dependency.  Editorial formatting only, in addition: an AMS \texttt{amsart} preamble that restates the article's own declarations and macros used by the retained lines, namely the environments \texttt{proposition}, \texttt{lemma} on separate counters, together with the standard packages those lines need.  The retained environments print in this document as Proposition 1, Lemma 1 in order of appearance, whereas the article prints the same items with the numbering of its own full document.  The complete native source of the article as distributed in the arXiv e-print for this exact version is bundled unchanged as \texttt{sources/176245/source0.tex}.
\begin{proposition}\label{prop:concentration-degrees}
	There exist two positive constants $C_1$ and $C_2$ such that, \whp,
	\begin{equation}\nonumber
		\max_{i \le m}\max_{x\in V_i}\max\{D_x^+,D_x^-,D_{x,i}^-\} \le C_2\lambda\log(n)\quad\text{and}\quad \min_{i \le m}\min_{x\in V_i}\min\{D_x^+,D_x^-,D_{x,i}^-\} \ge C_1\lambda\log(n)\,.
	\end{equation}
\end{proposition}

	\begin{lemma}\label{lemma:tree-like}
	For a digraph $S=(V,E)$, let $\trex(S):= |E|-|V|+1$ denote the tree-excess of $S$ and, for $\eps>0$,
	\begin{equation} \label{eq:G+}  
			\mathcal{G}_\eps:=\{\forall x \in V,\trex(\mathcal{B}^+_x({4\eps\tent}))< 2 \}.
	\end{equation}
	Let $(X_t)_{t \in \N}$ denote the SRW on $G_i$ or $G$ and, for $t \le {2\eps\tent}$, consider the set and related event
	\begin{align} \label{eq:Veps}
		\veps \coloneqq \left\{ x \in V \mid \trex(\mathcal{B}^+_{x}({4\eps\tent}))=0 \right\}
		, \qquad
		\mathcal{V}_{\eps}\coloneqq \{\max_{x \in V} \quenchG_x(X_t \notin \veps)\le 2^{-t}\}\,.
	\end{align}
	For a sufficiently small $\eps>0$ (independent of $n$), it holds $\p\left(\mathcal{G}_\eps \cap  \mathcal{V}_{\eps}\right) = 1-o(1)$.
\end{lemma}

\begin{lemma} \label{lemma:annealed-law}
For any $x\in V$, if $t \ll \sqrt{n}$,
\begin{equation}
	\pan_x(X_t=z,\mathcal{C}_t)=\frac{\semi^t(c(x),c(z))}{n}\left(1+O\left(\tfrac{t^2}{n}\right)\right)\,,\qquad \forall z\in V\,,
\end{equation}
where 
\begin{equation} \label{eq:Q-def}
	\semi^ t(i,j) = \frac{1+(m\one_{i=j}-1)(1-\frac{m}{m-1}\alpha)^{t}}{m}, \qquad i,j \le m, \quad t\ge0,
\end{equation}

\end{lemma}

\begin{equation} 
	\semi^ t(i,j) = \frac{1+(m\one_{i=j}-1)(1-\frac{m}{m-1}\alpha)^{t}}{m}, \qquad i,j \le m, \quad t\ge0,
\end{equation}

    \begin{proposition} \label{prop:quenched-law}
		For any $\alpha=\alpha_n\in[0,1]$, let $1\ll t \ll \sqrt{n}\log(n)^{-2} $. Then, 
		\begin{equation}\label{eq:315}
		\max_{i \le m}\max_{x \in V}
		\left| \quenchG_x(X_t \in V_i) - \semi^t(c(x),i)\right|
		=\op(1)\,,
		\end{equation}
		where, $\semi^t(c(x),i)$ is defined as in Eq.~\eqref{eq:Q-def}.
	\end{proposition}

		\begin{proof} 
		Notice that, since $m\asymp 1$, it suffices to prove \eqref{eq:315} for a fixed $i\le m$.
            We first consider the case in which $\alpha\ll 1$. Let $K=\lfloor \log^2 n \rfloor$ and call $\cE$ the event that all the vertices in the graph have out-degree at least $C_1\log(n)$ and $\cD$ the event in which each vertex has a tree-excess at most $1$ in its out-neighborhood of height $4$. Recall that, thanks to  Proposition \ref{prop:concentration-degrees} and Lemma \ref{lemma:tree-like}, $\PP(\cD\cap \cE)=1-o(1)$.
			Then
			\begin{equation}
			\p(\one_{\cE\cap\cD}\quenchG_x(X_t\in V_i)\ge \semi^t(c(x),i)+ \delta) \le \frac{\E[\one_{\cE\cap\cD}\quenchG_x(X_t\in V_i)^K]}{(\semi^t(c(x),i)+ \delta)^K}.
			\end{equation}
			Consider now $K$ annealed random walks $(X^{(\ell)})_{\ell\le K}$ starting at some $x\in V$. We let the walks evolve---all starting from $x$---one after the other, for a time $t$. Clearly, these walks are not independent, but each is independent of the previous ones \emph{conditionally on the environment discovered so far}.
			Let us define the family of events $(B_\ell)_{\ell \le K}$. For each $\ell \le K$, $B_\ell$ is the event that:
			\begin{enumerate}
				\item[(i)] $X^{(j)}_t \in V_i$ for each $j \le \ell$;
				\item[(ii)] all the out-degrees of the vertices visited by the first $\ell$ walks are at least $C_1\log(n)$ (cf.~Prop.~\ref{prop:concentration-degrees});
				\item[(iii)] the out-neighborhood of $x$ of height $4$ discovered by the first $\ell$ trajectories has a tree excess at most $1$.
			\end{enumerate}
			With this definitions we have
			\begin{equation}
			\E[\one_{\cE\cap\cD}\quenchG_x(X_t\in V_i)^K] \le \pan_x(B_K)=\pan_x(B_1)\prod_{{\ell}=2}^{K} \pan_x( B_{{\ell}}| B_{{\ell}-1})\,.
			\end{equation}
			We start by noting that, given $B_{{\ell}-1}$, the event $B_{{\ell}}$ is contained in the union of the following events:
			\begin{enumerate}
				\item[(1)] For all times $s\le 4$ the $\ell$-th trajectory is always in a vertex already visited by one of the previous walks.
                Since we are working on the events (ii) and (iii), for any couple $(x,y)$ there are at most $2$ paths of length $4$ joining $x$ to $y$ and each such path has a weight $\le (C_1\log(n))^{-4}$. The probability of the event described above is thus upper bounded by
            %    since, under the event in (iii), for any couple $(x,y)$ there are at most $2$ paths of length $4$ joining $x$ to $y$, and we are under the event in (ii),
            %this happens with probability at most 
        $$8K(C_1\log(n))^{-4}=o(1)\,.$$
				\item[(2)] The event in (1) does not occur, i.e., there exists some $s\le 4$ such that the walk visits an unvisited vertex, and there exists a time $s'\in(s,t]$ at which the trajectory intersects again one of the previous walks (including itself); this happens with probability less than 
                $$\frac{Kt^2}{n}=o(1)\,.$$
                \item[(3)] None of the events above occurs, and at time $s=4$ the walk is out of $V_{c(x)}$.  
                Since  $\alpha\ll 1$, this happens with probability at most
                $$1-(1-\alpha)^4=o(1)\,.$$
				\item[(4)] None of the event above occurs, yet at time $t$ the $\ell$-th trajectory is found in $V_i$; thanks to Lemma \ref{lemma:annealed-law}, this happens with probability at most
				\begin{align*}
					\max_{s\le 4}\max_{y\in V_{c(x)}}\pan_y(X_{t-s}\in V_i,\mathcal{C}_{t-s})&=\max_{s\le 4}\semi^{t-s}(c(x),i)\left(1+O\left(\frac{t^2}n\right)\right)\\
				&=\max_{\s\le 4}\semi^{t-s}(c(x),i)+o(1)\\
					&=\semi^{t}(c(x),i)+o(1)\,,
				\end{align*}
                where the first identity comes from the fact that, by symmetry, the annealed probability on the \lhs is independent of $y\in V_{c(x)}$; the second one follows from the fact that $\semi^t(c(x),i)\le 1$ and $t^2\ll n$;
                the third one uses that $t\gg 1$. 
			\end{enumerate}
			In conclusion,  uniformly over ${\ell}\le K$ we have $\pan_x(B_{\ell}|B_{{\ell}-1})=\semi^t(c(x),i)+o(1)$.
			Thanks to the choice of $K$, we conclude that, being $\delta>0$ fixed,
			\begin{equation}
	\p(\one_{\cE\cap\cD}\quenchG_x(X_t\in V_i)\ge\semi^t(c(x),i)+ \delta)
			\le \left(\frac{\semi^t(c(x),i)+o(1)}{\semi^t(c(x),i)+\delta}\right)^K 
			=o(n^{-1}).
			\end{equation}
			Therefore
   			\begin{equation}\label{eq:319}
            \begin{split}
				&\p\big(\max_{x \in V}\quenchG_x(X_t\in V_i)\ge \semi^t(c(x),i)+\delta\big)
				\le \\
                &\qquad\qquad \qquad\p(\one_{\cE\cap\cD}\max_{x \in V}\quenchG_x(X_t\in V_i)\ge \semi^t(c(x),i)+ \delta)+\p(\cE\cup\cD)=o(1)\,.
                \end{split}
			\end{equation}
            			To prove a uniform lower bound on $\quenchG_x(X_t\in V_i)$, one can consider the events
            \begin{equation}
        \bar{\mathcal{E}}_{x,i,\delta}=\left\{ \quenchG_x(X_t\in V_i)\le\semi^t(c(x),i)- \delta\right\}
            \end{equation}
            and
\begin{equation}
\hat{\mathcal{E}}_{x,j,\delta}=\left\{\quenchG_x(X_t\in V_j)\ge\semi^t(c(x),j)+   \frac{\delta}{m-1}\right\}\,.
            \end{equation}
Clearly, for any $i\le m$, $x\in V_i$ and $\delta>0$,  $\cup_{j\neq i}\hat{\mathcal{E}}_{x,j,\delta}\supseteq\bar{\mathcal{E}}_{x,i,\delta} $.
Therefore,
\begin{equation}\label{eq:322}
    \begin{split}
        \p\big(\min_{x \in V}\quenchG_x(X_t\in V_i)\le \semi^t(c(x),i)-\delta\big)&= \p\big(\cup_{x\in V}\bar{\mathcal{E}}_{x,i,\delta}\big)\\
        &\le\p\big(\cup_{j\neq i}\cup_{x\in V}\hat{\mathcal{E}}_{x,j,\delta}\big)\\
        &\le m\max_{j\le m}\p\big(\cup_{x\in V}\hat{\mathcal{E}}_{x,j,\delta}\big)\,.
    \end{split}
\end{equation}
Since $m\asymp 1$ and since the probability on the \rhs of \eqref{eq:322} coincides, replacing $\delta$ with $\frac{\delta}{m-1}$, with the one on the \lhs of \eqref{eq:319}, we conclude that
\begin{equation}\label{eq:323}
    \begin{split}
        \p\big(\min_{x \in V}\quenchG_x(X_t\in V_i)\le \semi^t(c(x),i)-\delta\big)&= o(1)\,.
    \end{split}
\end{equation}
Therefore, in the case $\alpha\ll1$, \eqref{eq:315} follows from \eqref{eq:319} and \eqref{eq:323} and a union bound over $i\le m$.

			 We are left to consider the case $\alpha\asymp 1$. In this case, $Q^t(j,i)=\frac1m+o(1)$ for any $t\gg1$. We can argue as above, but rather than the events (1), (2), (3) and (4), we can consider the events (1), (2) and (4'), where
			 \begin{itemize}
			 	\item[(4')] The events (1) and (2) do not occur, yet at time $t$ the $j$-th trajectory is found in $V_i$; thanks to Lemma \ref{lemma:annealed-law}, this happens with probability at most
			 	\begin{align*}
			 		\max_{s\le 4}\max_{y\in V}\pan_y(X_{t}\in V_i,\mathcal{C}_{t})&=\max_{s\le 4}\max_{j \le m}\semi^{t-s}(j,i)\left(1+O\left(\frac{t^2}n\right)\right)\\
			 		&=\frac1m+o(1)\,.
			 	\end{align*}
			 \end{itemize}
			 This completes the proof.
		\end{proof}

\end{document}
